% =============================================================================
% CHAPTER 1 — INTRODUCTION
% =============================================================================
\chapter{Introduction}
\label{ch:intro}

\section{Introduction}

India's retail sector is anchored by approximately twelve million kirana stores — neighbourhood
grocery and general merchandise outlets that serve as the primary point of purchase for a large
proportion of the country's population. These stores collectively account for over ninety percent
of the food and grocery retail market by volume and represent a cornerstone of India's informal
economy. Despite their economic significance, the vast majority of these establishments continue
to operate without the benefit of any digital inventory management system. Inventory decisions are
typically made on the basis of the owner's intuition, physical counting, or entries in paper
ledgers, practices that are inherently error-prone and incapable of adapting to the dynamic
consumer demand patterns of today's retail environment.

The consequences of this operational gap are commercially significant. A stockout event — where
a product runs out before the next replenishment order arrives — results in a direct loss of
revenue and a deterioration of customer trust. Conversely, overstocking ties up working capital
in slow-moving goods and increases the risk of spoilage for perishable items. Without a
systematic approach to tracking sales velocity, predicting replenishment needs, and evaluating
supplier performance, kirana store owners are left managing these risks entirely through manual
estimation.

Enterprise Resource Planning (ERP) solutions that address these challenges for large retailers
are commercially available but are wholly unsuitable for the kirana store context. Monthly
subscription costs ranging from several thousand to tens of thousands of rupees, complex
onboarding processes, and hardware requirements such as barcode scanners and RFID readers place
these systems well beyond the reach of a typical small store owner. Simpler Point-of-Sale (POS)
applications record transactions but lack the predictive analytics and automated alerting
capabilities necessary to prevent stockout situations proactively.

The emergence of low-code automation platforms, cloud-hosted databases, and large language model
APIs accessible at low or zero cost has created a new opportunity: it is now technically feasible
to build a fully functional, automated inventory intelligence system for a fraction of the cost
of traditional enterprise software. StockSense AI is designed to exploit precisely this
opportunity.

\section{Problem Statement}

Small retail stores in India lack an affordable and automated inventory management system that can predict stockout dates, generate supplier-aware reorder suggestions, score vendors, and deliver alerts through channels like WhatsApp. Existing solutions are financially and operationally inaccessible, while basic tools only provide passive recording without intelligence.

The core problem to be solved is therefore: \textit{How can an automated, AI-augmented inventory
intelligence pipeline be built and deployed for kirana stores at an infrastructure cost below
\rupee{}500 per month, requiring no technical expertise to operate on a day-to-day basis?}

\section{Objectives}

The following specific objectives have been established to guide the design and development of
StockSense AI:

\begin{enumerate}[label=\textbf{\arabic*.}, leftmargin=*, itemsep=4pt]
  \item To design and implement an automated pipeline that calculates stockout dates for all
        inventory items using daily sales velocity and updates product risk classifications
        (HIGH, MEDIUM, LOW) without manual intervention.

  \item To develop an AI reorder recommendation engine using Groq Llama 3.1 8B that generates
        natural-language suggestions with supplier selection reasoning based on stock levels,
        demand patterns, and performance grades.

  \item To build a supplier evaluation module using the Weighted Sum Model (WSM) that assigns
        composite performance scores and letter grades to suppliers based on reliability,
        pricing, and delivery speed criteria.

  \item To implement a real-time WhatsApp alert delivery system using the Twilio API that
        notifies store owners of critical stock conditions immediately upon pipeline execution,
        incorporating a 24-hour anti-spam cooldown mechanism.

  \item To create AI vision-powered features — product onboarding via photograph and supplier
        invoice scanning using Google Gemini~2.5 Flash — that automate data entry and eliminate
        manual stock update procedures.

  \item To deploy the complete system at a total infrastructure cost below \rupee{}500 per
        month using free-tier and low-cost cloud services, thereby ensuring financial
        accessibility for small retail businesses.
\end{enumerate}

\section{Scope of the Project}

StockSense AI is scoped as a complete, production-ready prototype targeting single-store kirana
operations in India. The scope encompasses the following dimensions:

\textbf{Functional Scope:} The system covers the full inventory intelligence cycle — from data
ingestion and stockout prediction to risk classification, AI-driven reorder suggestion, supplier
scoring, and WhatsApp alert delivery. It includes a browser-based React frontend providing
dashboards, product management, supplier analytics, and system logs. Automated daily orchestration
of the entire pipeline is included.

\textbf{AI and Automation Scope:} The project incorporates two distinct AI modalities: a text
generation model (Groq Llama~3.1~8B) for reorder reasoning and a multimodal vision model
(Google Gemini~2.5 Flash) for image-based data extraction. Workflow automation covers ten
modular n8n workflows, each responsible for a distinct stage of the intelligence pipeline.

\textbf{Data Scope:} The system manages eight database tables covering store profiles, products,
suppliers, stock alerts, reorder suggestions, stock transactions, daily snapshots, and system
logs. All tables are structured for multi-tenancy extension via \texttt{store\_id} foreign keys.

\textbf{Out of Scope:} The current implementation does not include a native mobile application,
multi-store tenancy with row-level security enforcement, demand forecasting using historical time
series models, or automated purchase order dispatch to suppliers. These are identified as
directions for future work.

\section{Organization of the Report}

The remainder of this report is organized as follows:

\textbf{Chapter 2 — Literature Survey} reviews published work on inventory management systems,
demand forecasting, large language model applications in procurement, IoT-based retail systems,
and low-code automation platforms. A comparative analysis table and identification of the research
gap are also presented.

\textbf{Chapter 3 — System Design} describes the four-layer architecture of StockSense AI,
presents a use case diagram description, data flow description, entity-relationship model, module
breakdown, and the complete database schema with all eight tables.

\textbf{Chapter 4 — Implementation} covers the technology stack with justification, development
environment configuration, module-wise implementation details for each of the ten workflows and
the frontend application, and key code snippets with explanations.

\textbf{Chapter 5 — Results and Discussion} presents system testing methodology, a structured
test case table with inputs and expected versus actual outcomes, descriptions of key application
screens, and a performance evaluation covering execution time, cost, and health score validation.

\textbf{Conclusion and Future Work} summarises the achievements of the project, evaluates them
against the stated objectives, acknowledges limitations honestly, and proposes specific, realistic
directions for future enhancement.

\textbf{References} lists all cited sources in APA format, including journal papers, conference
proceedings, and technical resources. \textbf{Appendices} provide key source code listings and
screenshot placeholders for all major application screens.
